\documentclass[journal, a4paper]{IEEEtran}

\usepackage[T1]{fontenc}
\usepackage[utf8]{inputenc}
\usepackage[spanish, es-noshorthands]{babel}
\usepackage{amsmath, amssymb, bm}
\usepackage{graphicx}
\usepackage{booktabs}
\usepackage{tabularx}
\usepackage{ragged2e}
\usepackage[american]{circuitikz}
\usepackage{pgfplots}
\pgfplotsset{compat=1.18}
\usepackage{hyperref}

\renewcommand{\tablename}{Tabla}

\begin{document}

\title{Análisis Experimental del Efecto de Carga y Máxima Transferencia de Potencia en una Red Equivalente de Thévenin}

\author{
    [Nombre del Estudiante 1], [Nombre del Estudiante 2]\\
    Ingeniería Mecatrónica, Universidad Católica Boliviana ``San Pablo'' -- Sede Tarija\\
    IMT-121 --- Circuitos Electrónicos I\\
    Fecha: \today \quad | \quad Grupo: [ID del Grupo]
}

\maketitle

\begin{abstract}
% Escriba un párrafo (80-120 palabras) indicando el barrido de resistencias realizado, la comparación entre SPICE y el banco, y la confirmación empírica del punto de máxima transferencia de potencia.
[Escribir resumen aquí]
\end{abstract}

\begin{IEEEkeywords}
Máxima transferencia de potencia, efecto de carga, eficiencia, equivalente de Thévenin, DC Sweep, LTSpice.
\end{IEEEkeywords}

\section{Introducción y Objetivo}
% Defina la dependencia de la potencia con respecto a la carga y plantee el objetivo del experimento.
[Escribir introducción aquí]

\section{Modelo del Circuito y Predicciones Analíticas}

\subsection{Definición del Circuito y Referencias}
% Incluya el equivalente de Thévenin conectado a una resistencia de carga variable (R_L). Indique polaridades y corrientes.
[Insertar esquemático de equivalente y carga]

\subsection{Relaciones Fundamentales}
La respuesta en carga de la red está gobernada por:
\begin{align}
    I_L &= \frac{V_{th}}{R_{th} + R_L} \\
    P_L(R_L) &= V_{th}^2 \frac{R_L}{(R_{th} + R_L)^2}
\end{align}
Matemáticamente, derivando respecto a $R_L$, la potencia se maximiza en la condición adaptada:
\begin{equation}
    R_{L,\text{max}} = R_{th}
\end{equation}

\subsection{Predicciones Analíticas}
El punto óptimo esperado para la red caracterizada es:
\begin{itemize}
    \item $R_{L,\text{max}} = R_{th} = $ [Valor]
    \item $P_{L,\text{max}} = $ [Valor]
\end{itemize}

\section{Metodología}
\subsection{Simulación (DC Sweep)}
% Describa el uso de LTSpice y el comando parametrizado (ej. .step param) para barrer múltiples valores de R_L y obtener la curva teórica.
[Descripción de simulación]

\subsection{Montaje y Procedimiento Experimental}
% Detalle cómo se realizó el barrido discreto de cargas físicas en la protoboard. Qué se mantuvo fijo (red) y qué variables (V_L, I_L) se registraron para calcular la potencia.
[Procedimiento del barrido físico]

\section{Resultados}

\begin{table*}[htbp]
    \centering
    \caption{Barrido Discreto: Evolución de Tensión, Corriente y Potencia}
    \label{tab:resultados_b}
    \renewcommand{\arraystretch}{1.2}
    \begin{tabularx}{\textwidth}{>{\centering\arraybackslash}p{2.5cm} *{5}{>{\centering\arraybackslash}X}}
        \toprule
        $\bm{R_L}$ \textbf{Medido} ($\Omega$) & $\bm{V_L}$ \textbf{Teórico} (V) & $\bm{V_L}$ \textbf{Medido} (V) & $\bm{I_L}$ \textbf{Calculado} (mA) & $\bm{P_L}$ \textbf{Teórico} (mW) & $\bm{P_L}$ \textbf{Medido} (mW) \\
        \midrule
        $R_{L1}$ & & & & & \\
        $R_{L2}$ & & & & & \\
        $R_{L3} \approx R_{th}$ & & & & & \\
        $R_{L4}$ & & & & & \\
        $R_{L5}$ & & & & & \\
        \bottomrule
    \end{tabularx}
\end{table*}

\begin{figure}[htbp]
    \centering
    % El estudiante debe reemplazar este entorno por su gráfico (P_L vs R_L) generado desde Python, MATLAB o Excel, mostrando los puntos medidos sobre la curva teórica/simulada. Priorice la legibilidad de ejes y leyenda.
    \begin{tikzpicture}
        \begin{axis}[
            width=\columnwidth,
            height=6cm,
            xlabel={Resistencia de Carga $R_L$ ($\Omega$)},
            ylabel={Potencia $P_L$ (mW)},
            grid=major,
            legend pos=north east
        ]
        % \addplot table {data.txt}; % Placeholder para datos reales
        \end{axis}
    \end{tikzpicture}
    \caption{Curva de Potencia en función de la resistencia de carga $R_L$. Se indican el máximo predicho y los puntos discretos medidos experimentalmente.}
    \label{fig:curva_potencia}
\end{figure}

\section{Discusión}
% Responda explícitamente:
% 1. ¿Coincidió razonablemente el pico experimental de potencia con el valor medido de R_th?
% 2. Analice el efecto de carga: ¿Qué le sucede a V_L y a la potencia útil cuando R_L es excesivamente grande o excesivamente pequeña respecto a R_th?
% 3. ¿Cómo se comportó la eficiencia de transferencia en el punto exacto de máxima potencia?
[Escribir discusión aquí]

\section{Conclusiones}
% Escriba de 3 a 5 conclusiones técnicas precisas respaldadas por los datos.
[Escribir conclusiones aquí]

\begin{thebibliography}{00}
\bibitem{guia_b} Instructor del Curso, \emph{Experiencia Integrada B}, Universidad Católica Boliviana, 2026.
\end{thebibliography}

\end{document}
